% BOOK-FILE-ID: f-daeb39e05261
% section: 発展的な解法

\section{発展的な解法}\label{節：発展的な解法}
ここからは発展的な漸化式を扱う.
係数が$n$によって変わる式,ガウス記号を含む式,関係する項が多い式などである.
難しく見えるかもしれないが,やることは変わらない.既知の形をつくるのだ.
これまでの解法をどう組み合わせれば既知の形に帰着できるか,一つ一つ考えていこう.\footnote{ここで扱う\bookterm{recurrence}{漸化式}は,一部の難関大でも出題される.
  基本的に誘導がつくので,その誘導が何を既知の形にしようとしているのかを読み取ることが大切である.}
\BookImport{f-847298b7c959}
\BookImport{f-0497b0842775}
\BookImport{f-5d9116720103}
\BookImport{f-ff87377cd31b}
\BookImport{f-59858680990d}
